\section{Key Derivation and Deterministic Nonce}
\label{sec:kdf-nonce}

\subsection{Key schedule}

Let \code{ss\_mlkem} be the 32-byte shared secret returned by ML-KEM-1024 \code{Decaps} (or \code{Encaps} on the sender side). Compute:
\begin{align*}
\mathrm{salt} &= \mathrm{SHAKE256}(\text{``pqcrypt v1 header''} \,\|\, \mathrm{header\_bytes}[0\!:\!16],\ 32) \\
\mathrm{prk}  &= \mathrm{HKDF\text{-}Extract}(\mathrm{salt},\ \texttt{ss\_mlkem}) \\
K_{\mathrm{aead}}   &= \mathrm{HKDF\text{-}Expand}(\mathrm{prk},\ \text{``pqcrypt v1 aead k''},\ 32) \\
\mathrm{NONCE}      &= \mathrm{HKDF\text{-}Expand}(\mathrm{prk},\ \text{``pqcrypt v1 aead n''},\ 12) \\
K_{\mathrm{commit}} &= \mathrm{HKDF\text{-}Expand}(\mathrm{prk},\ \text{``pqcrypt v1 commit''},\ 32)
\end{align*}
The PRF is HMAC-SHA3-256 \citep{fips202}.

\subsection{Salt binding}

The salt mixes a 16-byte project-and-version constant with the full envelope header (timestamp, variant, suite, counter, KEM CT length, AD length). Any header bit-flip therefore invalidates $K_{\mathrm{aead}}$ before the AEAD tag is even checked, providing a second line of defense against header tampering. The constant prevents cross-protocol reuse of the KDF.

\subsection{Why deterministic nonce}

Each ML-KEM-1024 encapsulation produces a fresh ephemeral shared secret, so $K_{\mathrm{aead}}$ is unique per envelope. Deterministic derivation of \code{NONCE} from \code{prk} therefore cannot cause nonce reuse under the same key. This:
\begin{itemize}
\item Removes the encrypt-side CSPRNG dependency beyond ML-KEM's internal \code{Encaps} randomness.
\item Matches PASETO v4.local's philosophy \citep{paseto-v4}: deterministic derivation under a per-message fresh key is structurally safer than per-message random nonce generation.
\item Eliminates a class of misuse: an implementer who accidentally reuses a CSPRNG seed cannot create a nonce-collision oracle, because the seed cannot survive ML-KEM \code{Encaps} \citep{fips203}.
\end{itemize}

\subsection{Multi-user / multi-key bounds}

The single-key birthday bound on AES-256-GCM is $2^{64}$ blocks; PQCrypt's plaintext cap of $2^{32}$ bytes per envelope keeps each invocation well below the bound. The multi-user attack analysis of Bose, Hoang, and Tessaro \citep{bose2018muser,hoang2024gcmsst} concerns adversaries who attack many keys simultaneously; PQCrypt's per-envelope fresh $K_{\mathrm{aead}}$ removes the long-key-lifetime axis of that attack class.
